% Condensed from numerical/tables/ising_large_parameters.tex; numerical values copied without refitting.
\begin{table}[htbp]\centering\footnotesize
\setlength{\tabcolsep}{3.5pt}
\begin{tabular}{lrrrrrrr}\toprule
Direction & $\widehat\beta$ & $C_{\rm free}$ & $c_0^{\rm CFT}$ & $c_0(F_1)$ & $c_0(F_2)$ & $c_1(F_2)$ & $\eta_0(F_2)$\\\midrule
$I\mid \sigma$ & 0.26465 & 0.71994 & 0.61965 & 0.65339 & 0.61463 & 0.20243 & -0.81098\\
$\sigma\mid I$ & 0.26465 & 0.71995 & 0.61965 & 0.65339 & 0.61462 & 0.20245 & -0.81139\\
$I\mid \varepsilon$ & 1.9985 & 3.2548 & 3.2899 & 3.2861 & 3.289 & -1070.2 & -0.024918\\
$\varepsilon\mid I$ & 2.0013 & 3.3242 & 3.2899 & 3.2964 & 3.2938 & 959.21 & 0.11818\\
$\sigma\mid \varepsilon$ & 0.25058 & 0.15691 & 0.15491 & 0.15631 & 0.15595 & 0.0018716 & 0.67037\\
$\varepsilon\mid \sigma$ & 0.25544 & 0.16288 & 0.15491 & 0.15712 & 0.15365 & 0.018087 & -0.81285\\
\bottomrule\end{tabular}
\caption{Ising large-interval fits on $0<\ep\leq0.002$, with 260 geometries per direction. The free fit is $C_{\rm free}\ep^{\widehat\beta}$; $F_1,F_2$ impose the powers in Table~\ref{tab:num-powers}. All amplitudes include powers of $\pi$. The final column is the signed leading-amplitude discrepancy of $F_2$ in percent, Eq.~\eqref{eq:num-residuals}. The $c_1$ column gives a fitted next amplitude, not a theoretical prediction.}
\label{tab:num-ising-large}
\end{table}
